\documentclass[11pt,a4paper]{article}
\usepackage[utf8]{inputenc}
\usepackage{amsmath,amssymb,amsthm}
\usepackage{listings}
\usepackage{geometry}
\usepackage{hyperref}
\geometry{margin=2.5cm}

\newtheorem{theorem}{Theorem}
\newtheorem{lemma}[theorem]{Lemma}

\lstset{
  language=C++,
  basicstyle=\ttfamily\small,
  keywordstyle=\bfseries,
  breaklines=true,
  frame=single,
  numbers=left,
  numberstyle=\tiny,
  tabsize=4,
  showstringspaces=false
}

\title{IOI 2008 -- Pyramid Base}
\author{Solution}
\date{}

\begin{document}
\maketitle

\section{Problem Statement}
Given an $R \times C$ grid with $P$ rectangular obstacles (each with a removal cost) and a budget $B$, find the largest side length $s$ of an axis-aligned square that can be placed on the grid such that the total cost of overlapping obstacles is at most $B$.

\section{Solution}

\subsection{Binary Search on Side Length}
\begin{lemma}[Monotonicity]
If a square of side $s$ can be placed with cost $\le B$, then any square of side $s' < s$ can be placed at the same position with cost $\le B$ (it overlaps a subset of the same obstacles). This justifies binary search on $s$.
\end{lemma}

\subsection{Feasibility Check via Sweep Line}
For a given side length $s$, determine if any placement has cost $\le B$:
\begin{enumerate}
    \item For each obstacle $i$ with bounding box $[x_1^i, y_1^i] \times [x_2^i, y_2^i]$ and cost $c_i$, the set of top-left corners $(r, c)$ whose $s \times s$ square overlaps this obstacle forms a rectangle. Specifically:
    \[
    r \in [\max(1, x_1^i - s + 1),\; \min(R - s + 1, x_2^i)], \quad
    c \in [\max(1, y_1^i - s + 1),\; \min(C - s + 1, y_2^i)].
    \]
    \item Create sweep-line events: at row $r_{\text{start}}$, add $c_i$ to the column range; at row $r_{\text{end}}+1$, subtract $c_i$.
    \item Sweep rows top to bottom, maintaining a \textbf{segment tree} supporting range-add and global-minimum queries. If the global minimum $\le B$ at any row, side length $s$ is feasible.
\end{enumerate}

\subsection{Complexity}
\begin{itemize}
    \item \textbf{Time:} $O\bigl(\log(\min(R,C)) \cdot (R + P) \cdot \log C\bigr)$.
    \item \textbf{Space:} $O(R + C + P)$.
\end{itemize}

\section{C++ Solution}

\begin{lstlisting}
#include <bits/stdc++.h>
using namespace std;

struct SegTree {
    int n;
    vector<long long> mn, lazy;

    SegTree(int n) : n(n), mn(4*n, 0), lazy(4*n, 0) {}

    void push(int v) {
        if (lazy[v]) {
            for (int c : {2*v, 2*v+1}) {
                mn[c] += lazy[v];
                lazy[c] += lazy[v];
            }
            lazy[v] = 0;
        }
    }

    void update(int v, int l, int r, int ql, int qr, long long val) {
        if (ql > r || qr < l) return;
        if (ql <= l && r <= qr) {
            mn[v] += val; lazy[v] += val; return;
        }
        push(v);
        int mid = (l + r) / 2;
        update(2*v, l, mid, ql, qr, val);
        update(2*v+1, mid+1, r, ql, qr, val);
        mn[v] = min(mn[2*v], mn[2*v+1]);
    }

    void update(int ql, int qr, long long val) {
        if (ql > qr) return;
        update(1, 1, n, ql, qr, val);
    }

    long long query() { return mn[1]; }
};

int main(){
    ios::sync_with_stdio(false);
    cin.tie(nullptr);

    int R, C, P;
    long long B;
    cin >> R >> C >> P >> B;

    vector<int> x1(P), y1(P), x2(P), y2(P);
    vector<long long> cost(P);
    for (int i = 0; i < P; i++)
        cin >> x1[i] >> y1[i] >> x2[i] >> y2[i] >> cost[i];

    int lo = 1, hi = min(R, C), ans = 0;
    while (lo <= hi) {
        int s = (lo + hi) / 2;
        int maxR = R - s + 1, maxC = C - s + 1;
        if (maxR < 1 || maxC < 1) { hi = s - 1; continue; }

        // Create events: events[row] = list of (col_start, col_end, delta)
        vector<vector<tuple<int,int,long long>>> events(maxR + 2);
        for (int i = 0; i < P; i++) {
            int rs = max(1, x1[i] - s + 1);
            int re = min(maxR, x2[i]);
            if (rs > re) continue;
            int cs = max(1, y1[i] - s + 1);
            int ce = min(maxC, y2[i]);
            if (cs > ce) continue;
            events[rs].emplace_back(cs, ce, cost[i]);
            if (re + 1 <= maxR)
                events[re + 1].emplace_back(cs, ce, -cost[i]);
        }

        SegTree seg(maxC);
        bool feasible = false;
        for (int r = 1; r <= maxR && !feasible; r++) {
            for (auto& [cs, ce, delta] : events[r])
                seg.update(cs, ce, delta);
            if (seg.query() <= B) feasible = true;
        }

        if (feasible) { ans = s; lo = s + 1; }
        else hi = s - 1;
    }

    cout << ans << "\n";
    return 0;
}
\end{lstlisting}

\section{Notes}
The segment tree supports lazy range-add and global-minimum queries. Each obstacle creates two sweep-line events (entry and exit). The total work per binary-search iteration is $O((R + P) \log C)$, giving an efficient overall solution.

For $B = 0$, the problem reduces to finding the largest obstacle-free square, which this approach handles as a special case.

\end{document}
